\documentclass[10pt,a4paper]{amsart}
\usepackage[margin=19mm]{geometry}
\usepackage{amsmath,amssymb,mathtools}
\usepackage[scr=rsfs,cal=euler]{mathalfa}
\usepackage{xfrac}
\usepackage[unicode,hidelinks,bookmarks=false]{hyperref}
\newcommand{\ie}{i.e.\ }
\newcommand{\cf}{cf.\ }
\newcommand{\resp}{respectively\ }
\newcommand{\Subsection}{\subsection}
\providecommand{\nobreakdash}{\nobreak\hbox{-}}
\setlength{\emergencystretch}{2em}
\multlinegap0pt
\usepackage{amssymb}
\usepackage[normalem]{ulem}
\newtheorem{proposition}{Proposition}
\newtheorem{theorem}{Theorem}
\DeclareMathOperator{\GL}{GL}
\newcommand{\NN}{\mathbb{N}}
\newcommand{\ZZ}{\mathbb{Z}}
\title{Lecture hall polytopes and\\Lakshmibai-Seshadri paths}
\author{Rocco Chiriv\`i, Martina Costa Cesari, Xin Fang, Peter Littelmann}
\date{Source: arXiv:2609.03195v1 (CC BY 4.0)}
\begin{document}
\maketitle

\noindent\emph{Source and licence.} Lecture hall polytopes and\\Lakshmibai-Seshadri paths. Version arXiv:2609.03195v1 (CC BY 4.0), distributed by arXiv under the Creative Commons Attribution 4.0 International licence, http://creativecommons.org/licenses/by/4.0/, as recorded in the arXiv licence metadata for this exact version and shown on the abstract page https://arxiv.org/abs/2609.03195v1. Full licence notice: \texttt{licenses/arxiv-2609.03195-CC-BY-4.0.txt}.\par\smallskip

\noindent\textit{Editorial scope.} Counted unit: the article's theorem thm:invariant-ring-Gamma (source0.tex lines 399--401). The article's complete original proof of it is delivered with it (source0.tex lines 402--423, one \texttt{proof} environment of 22 lines). No mathematics is written, repaired or re-derived: the statement and the proof are the article's own lines, in the article's own notation, and no retained line is substituted or re-flowed.  Retained uncounted as the source-local context the proof needs: lines 228--241.  Nothing was omitted from any retained span, so the package carries no omission record.  No retained line contains a citation, so the wrapper carries no bibliography and no bibliography entry.  No retained line contains a figure environment, an image-inclusion command or drawing code, so no image is delivered and the package declares no asset dependency.  Editorial formatting only, in addition: an AMS \texttt{amsart} preamble that restates the article's own declarations and macros used by the retained lines, namely the environments \texttt{proposition}, \texttt{theorem} on separate counters, together with the standard packages those lines need.  The retained environments print in this document as Proposition 1, Theorem 1 in order of appearance, whereas the article prints the same items with the numbering of its own full document.  The complete native source of the article as distributed in the arXiv e-print for this exact version is bundled unchanged as \texttt{sources/209276/source0.tex}.
\begin{proposition}\label{prop:LS-characterization}
    A sequence of non-negative real numbers $\pi = (a_1,\dots,a_{n+1})$ is an L-S path if and only if
    \[
    \left\{
    \begin{array}{rcl}
        b_1 a_1 & \in &  \NN,\\
        b_2(a_1 + a_2) & \in &  \NN,\\
        & \vdots\\
        b_n(a_1 + \cdots + a_n) & \in &  \NN,\\
        a_1 + \cdots + a_{n+1} & \in &  \NN.
    \end{array}
    \right.
    \]
\end{proposition}

\begin{theorem}\label{thm:invariant-ring-Gamma}
The algebra $A(b)$ is the subalgebra of $\Gamma$--invariants in $ R$, i.e. $A(b) =  R^{\Gamma}$. Moreover, the subgroup $\Gamma \subset \GL( V)$ is free of pseudo-reflections.
\end{theorem}

\begin{proof}
The action of $\Gamma$ is diagonal on the monomial basis of $R$, so a general element of $R$ is $\Gamma$--invariant if and only if each monomial appearing in it is $\Gamma$--invariant. Hence it is enough to check which monomials $t^\pi = t_1^{a_1}\cdots t_{n+1}^{a_{n+1}}$, $\pi = (a_1, \ldots,a_{n+1})$, are $\Gamma$--invariant.
We have, $t_k^{a_k} = (t_k^{1/N_k})^{m_k}$ with $m_k := N_k a_k \in \ZZ_{\ge 0}$, so
\[
\gamma_j(t^\pi) = \gamma_j((t_1^{1/N_1})^{m_1})\cdots \gamma_j((t_{n+1}^{1/N_{n+1}})^{m_{n+1}}) = \zeta^{b_j(m_1N/N_1 + \cdots + m_jN/N_j)}\, t^\pi.
\]
Hence $t^\pi$ is $\Gamma$--invariant if and only if
\[
b_j(m_1N/N_1 + \cdots + m_jN/N_j)
\]
is a multiple of $N$ for all $j = 1, \ldots, n+1$. But $m_kN/N_k = N a_k$ for all $k$, so the previous condition is equivalent to $b_j(a_1 + \cdots + a_j) \in \ZZ$ for all $j = 1, \ldots, n+1$. By Proposition~\ref{prop:LS-characterization}, this is exactly the condition for $\pi$ to be an L-S path. This completes the proof that $A(b) =  R^{\Gamma}$.

Now we show that $\Gamma$ is free of pseudo-reflections. For an arbitrary element $\gamma = \gamma_1^{\ell_1}\cdots \gamma_{n+1}^{\ell_{n+1}}$ of $\Gamma$
\[
\gamma(t_k^{1/N_k}) = \zeta^{N c_k /N_k}\, t_k^{1/N_k},\quad c_k := b_k \ell_k + \cdots + b_{n+1} \ell_{n+1}.
\]
So suppose that $\gamma$ has all eigenvalues $\lambda_1, \ldots, \lambda_{n+1}$ equal to $1$ but at most one of them; i.e. there exists $1\leq j\leq n+1$ such that $\lambda_k = 1$ for all $k\neq j$. We will show that $\gamma$ is the identity.

The condition $\lambda_k = 1$ is equivalent to $N_k\,|\,c_k$ and, since $N_k = \mathrm{lcm}(b_k,b_{k-1})$, this is equivalent to $b_k\,|\,c_k$ and $b_{k - 1}\,|\,c_k$. Furthermore, since $c_k = b_k\ell_k + \cdots + b_{n+1}\ell_{n+1}$, the condition $b_k\,|\,c_k$ is equivalent to $b_k\,|\,c_{k+1}$ (where we put $c_{n+2} := 0$).

So, denoting by ${\mathsf C}_k$ the condition $b_k\,|\,c_{k+1}$, we have proved that $\lambda_k = 1$ if and only if ${\mathsf C}_{k-1}$ and ${\mathsf C}_k$ hold. The condition ${\mathsf C}_{n+1}$ is always true since $b_{n+1} = 1$. So, our hypothesis on $\gamma$ implies that ${\mathsf C}_k$ and ${\mathsf C}_{k-1}$ hold for all $k\neq j$. Hence ${\mathsf C}_k$ holds for each $k=1,\ldots,n+1$, and $\gamma$ is the identity.
\end{proof}

\end{document}
